\section{Realted Work}\label{sec:related_work}

MLLMs have achieved impressive performance in visual understanding tasks by utilizing VLMs to process an image into a sequence of visual tokens, with lengths ranging from 64 to 729~\cite{radford2021learning,zhai2023sigmoid}, which are then input into LLMs~\cite{liu2024llavanext,yang2024qwen2,li2024llava,yao2024minicpm,chen2024far}.
However, in video understanding, a video contains hundreds or even thousands of frames, resulting in a long sequence of visual tokens.
As a result, feeding an entire video into Video-LLMs becomes impractical. 
Some research has explored extending the context length of LLMs~\cite{wan2023efficient,zhang2024longva,xiong2024effective,xue2024longvila,song2024moviechat},
but these methods still struggle with computational burden and the ``lost-in-the-middle'' issue, where important context may be overlooked in long sequences~\cite{liu2024st,yu2024frame}.



Most existing Video-LLMs rely on uniform frame sampling~\cite{chen2023internvl,li2023videochat,liu2024oryx,shu2024video}, 
which arbitrarily disregards the current query, often missing relevant frames and including irrelevant ones.
Traditional key-frame methods that detect significant transitions between adjacent frames also face this limitation~\cite{sheena2015key,nasreen2013key,wolf1996key}.
%
Some studies have attempted to use the query to retrieve the top-$k$ most relevant frames with VLMs but fail to consider frame diversity~\cite{wang2024weakly,yu2024self},
often resulting in redundant frame selection, and increased computational burden without meaningful information gain.
Some cluster-based methods attempt to regularize pair-wise diversity among frames but still fail to capture combinatorial relationships in a list-wise manner~\cite{keplerlab_katna,jin2024chat}. 
%
Recently, \citet{yu2024frame} proposed a frame selection method tailored for VILA~\cite{lin2024vila}, 
%Recently, \citet{yu2024frame, hu2025m} proposed frame selection methods tailored for specific Video-LLMs,
which requires collecting a supervised dataset and end-to-end training of a selector, limiting its generalizability to other Video-LLMs.
Furthermore, these method does not explicitly address list-wise diversity or sequentiality, both of which are essential for effective frame selection.



The DPP was introduced by~\citet{macchi1975coincidence} to model the distribution of fermions in thermal equilibrium, where no two fermions can occupy the same quantum state, resulting in an ``anti-bunching'' effect that can be interpreted as diversity and effectively modeled using DPP. 
DPP has been widely applied to characterized list-wise diversity in subset selection problems~\cite{kulesza2012determinantal,moss2023inducing,celis2018fair,li2024evaluation,chen2018fast}. 
However, standard DPP does not account for sequentiality. 
\citet{gong2014diverse} addressed this by segmenting the video and modeling local diversity within each segment conditioned on the selections from the preceding segment.
%
Nevertheless, allocating the total selection size \(k\) across segments remains challenging.
\citet{zheng2021k} treat it as an integer programming problem, relaxing problem constraints to find an approximate solution.
In contrast, we model DPP for sequentiality as an MDP, enabling optimal allocation of \(k\) through a dynamic programming algorithm with pseudo-polynomial time complexity. 
This approach avoids constraint relaxation and can find the optimal allocation efficiently in practice.



